% ============================================================================
% CostAware — IEEE conference-format report
% Compile on Overleaf (IEEEtran is built in) or locally:
%   pdflatex costaware_ieee && bibtex costaware_ieee && pdflatex costaware_ieee && pdflatex costaware_ieee
%
% Every number in Section V is produced by:
%   python scripts/trigger_coverage.py --seeds 50 --rate 0.25 --json results/tables/trigger_coverage_r025.json
% Sections VI–VII describe experiments that have NOT been run yet; do not add
% result numbers there until the harness (Phases 5–8) produces them.
% ============================================================================
\documentclass[conference]{IEEEtran}
\IEEEoverridecommandlockouts

\usepackage{cite}
\usepackage{amsmath,amssymb,amsfonts}
\usepackage{graphicx}
\usepackage{booktabs}
\usepackage{array}
\usepackage{makecell}
\usepackage{xcolor}
\usepackage{url}
\usepackage[hidelinks]{hyperref}

\newcommand{\yes}{\checkmark}
\newcommand{\no}{--}
\newcommand{\pt}{$\circ$}

\begin{document}

\title{Which Faults Does Adaptive Verification Notice?\\
A Matched-Fault Testbed for Cost-Aware Verification\\
in Long-Horizon Tool-Using LLM Agents}

% TODO: add co-authors (team members on SP-00-Origin/costAware) and supervisor.
\author{\IEEEauthorblockN{Krish Sharma}
\IEEEauthorblockA{\textit{Department of Computer Science and Engineering}\\
\textit{BML Munjal University}\\
Gurugram, India\\
krish.sharma.24cse@bmu.edu.in}
}

\maketitle

% ----------------------------------------------------------------------------
\begin{abstract}
Tool-using large language model (LLM) agents act on observations returned by external
tools, and a single corrupted observation can propagate into an irreversible
state-changing action. Verifying every step is expensive, so recent systems verify
\emph{selectively}, typically triggering on the risk tier of the tool being called or on
malformed tool output. We ask a question that prior selective-verification work does not
isolate: \emph{which kinds of tool faults do such triggers actually notice?} We present
CostAware, a deterministic testbed in which an agent completes long-horizon enterprise
procurement tasks through 17 risk-tiered tools while a seeded injector corrupts tool
outputs with eight fault types. Its key design property is \emph{matched faults}: faults
are keyed on the tool call itself, so the same agent action receives the same fault under
every verification policy, regardless of how many extra verification calls a policy makes.
A deterministic, LLM-free oracle judges the hidden true state. In an LLM-free trigger-coverage
analysis over 6{,}950 tool calls (20 tasks $\times$ 50 seeds), a risk-and-schema adaptive policy
verifies 43.7\% of steps and reaches 63.1\% of faulted steps, against 39.8\% for a fixed
every-third-step schedule at a similar 38.1\% rate. However, its content-based trigger fires on
0\% of \emph{silent} faults (well-formed but wrong values, stale records, or wrong statuses);
on these it does no better than the fixed schedule (42.1\% vs.\ 44.4\%) and is worse on two
fault types. We release the testbed and an evaluation protocol for comparing never-, always-,
scheduled-, and adaptively-verifying agents on reliability versus cost.
\end{abstract}

\begin{IEEEkeywords}
LLM agents, tool use, verification, fault injection, cost-aware inference, reliability
\end{IEEEkeywords}

% ----------------------------------------------------------------------------
\section{Introduction}

LLM agents increasingly execute multi-step workflows through tool calls: they query
records, compare options, and then commit state-changing actions such as creating or
releasing a purchase order. The agent's only view of the world is what the tools return.
If a tool returns a stale quote, a flipped vendor-suspension flag, or a mis-reported budget,
the agent can reason correctly and still act wrongly. Recent fault-injection studies
show that this is common: corrupted tool outputs substantially reduce task success
\cite{gupta2026reliabilitybench,wang2026agentnoisebench,agentchaos2026,zhu2026toolmaze}, a single
wrong parameter propagates to a wrong final answer most of the time \cite{gurram2026agentprop},
and agents notice explicit errors far more readily than plausible wrong values
\cite{kraishan2026loud}.

Verification---re-checking an observation or an action before relying on it---is the
obvious remedy, but it is not free. Each verifier call costs tokens and latency, and
enforcing checks on every step imposes a measurable ``verifier tax''
\cite{sah2026verifiertax}. A growing line of work therefore verifies selectively: gating
critique on action ambiguity \cite{park2026sag}, placing verifiers only on error-prone
workflow nodes \cite{ro2025sherlock}, re-checking only before backend writes
\cite{chen2026reconsider}, or escalating from a cheap deterministic check to an LLM judge
\cite{sieve2025,safeharness2026,verificationlayer2026}.

These triggers share an implicit assumption: that the signals they watch---the risk of the
tool being called, or whether its output is well formed---are correlated with the
observations that are actually wrong. This paper tests that assumption directly. Our
contributions are:

\begin{enumerate}
\item \textbf{A controlled procurement testbed} (Section~\ref{sec:testbed}): a deterministic
simulator with 17 risk-tiered tools, 20 tasks in three difficulty tiers, and an LLM-free
oracle. The simulator enforces workflow integrity but deliberately \emph{permits} policy
violations (e.g., ordering from a suspended vendor), so a misled agent can actually fail.
\item \textbf{A matched-fault protocol} (Section~\ref{sec:faults}): faults are derived from a
hash of the seed, tool, arguments, and call occurrence, so verification policies that make
different numbers of extra calls still face identical faults on identical agent actions,
and every episode can be replayed exactly.
\item \textbf{A trigger-coverage analysis} (Section~\ref{sec:coverage}) showing that
risk-tier and schema triggers are effectively blind to silent value faults: the content
trigger fires on none of them, and the adaptive policy covers them no better than a fixed
schedule.
\item \textbf{An evaluation protocol} (Section~\ref{sec:protocol}) for comparing verification
policies A--E on reliability versus cost, including signal ablations and a rate-matched
random-verification control.
\end{enumerate}

We do not yet report end-to-end success rates of LLM agents; Section~\ref{sec:protocol}
describes those experiments, and Section~\ref{sec:limits} states what the present results do
and do not show.

% ----------------------------------------------------------------------------
\section{Related Work}
\label{sec:related}

\textbf{Fault injection for tool-using agents.}
ReliabilityBench applies chaos-engineering faults (timeouts, rate limits, partial responses,
schema drift) and reports a reliability surface over repetition, perturbation, and fault
intensity \cite{gupta2026reliabilitybench}. AgentNoiseBench injects user- and tool-side noise
into existing benchmarks \cite{wang2026agentnoisebench}; AgentChaos injects faults at the LLM
API layer \cite{agentchaos2026}; ToolMaze distinguishes explicit/implicit and
transient/permanent tool failures and measures recovery cost \cite{zhu2026toolmaze};
AgentProp-Bench measures how an injected wrong parameter propagates \cite{gurram2026agentprop}.
Kraishan finds that agents respond to the error channel rather than to content, so
in-format faults pass unnoticed \cite{kraishan2026loud}, and SSCBench argues that fault-injection
results should report whether refuting evidence was available before first use
\cite{he2026sscbench}. These works \emph{measure} agent robustness; none compares verification
\emph{policies} under faults held fixed across policies, which is our focus. Our fault
taxonomy overlaps with theirs (timeouts, malformed output, wrong values) and we do not claim
it as a contribution.

\textbf{Selective verification and critique.}
SAG gates an external critic on action-level ambiguity (entropy and top-2 margin) and
improves success per token over no-critique and always-critique agents
\cite{park2026sag}; it is the closest work in spirit, but it targets embodied decision
benchmarks without tool-output corruption. Sherlock places verifiers on error-prone workflow
nodes identified by counterfactual fault analysis \cite{ro2025sherlock}. AVA allocates compute
across search, sampling, and verification cascades under a budget for reasoning tasks
\cite{patel2026ava}, and a communication-theoretic view maps retry, voting, and critique onto
classical reliability operators with a cost-aware router \cite{omidvar2026comm}. Earlier
self-correction work (Reflexion \cite{shinn2023reflexion}, Self-Refine
\cite{madaan2023selfrefine}, CRITIC \cite{gou2024critic}) verifies uniformly, and intrinsic
self-correction without external feedback is known to be unreliable \cite{huang2024cannot}.

\textbf{Guards on state-changing actions.}
Write-triggered reconsideration re-checks proposed backend writes on $\tau^2$-bench
\cite{chen2026reconsider}; SIEVE and SafeHarness escalate from deterministic checks to an LLM
adjudicator \cite{sieve2025,safeharness2026}; RACG gates which risk-tiered tools are exposed
at all \cite{racg2026}; and a verification wrapper with post-conditions and idempotency keys
addresses non-atomic failures such as timeouts after dispatch \cite{mansoor2026verifiedtools}.
Policy D in this paper is an instance of this family (risk-tier plus schema trigger), which is
precisely why measuring its blind spots is useful.

\textbf{Cost-aware cascades and environments.}
Cheap-then-expensive cascades originate in FrugalGPT \cite{chen2023frugalgpt} and have online,
regret-bounded variants \cite{omvv2026}. Enterprise and business environments for agents
include $\tau$-bench \cite{yao2024taubench}, EnterpriseBench \cite{enterprisebench}, and
Vending-Bench \cite{backlund2025vending}; ToolEmu emulates risky tool executions
\cite{ruan2024toolemu}.

Table~\ref{tab:positioning} summarizes the positioning. The individual ingredients exist; the
combination of matched faults across verification policies, a stateful environment where
faults can cause irreversible errors, and an LLM-free oracle is, to our knowledge, not
covered by prior work.

\begin{table}[t]
\caption{Positioning against the closest prior work (from published abstracts).
\yes: central; \pt: partial; \no: not addressed.}
\label{tab:positioning}
\centering
\setlength{\tabcolsep}{3pt}
\footnotesize
\begin{tabular}{@{}lcccccc@{}}
\toprule
Work & \makecell{Tool-\\output\\faults} & \makecell{State-\\changing\\actions} & \makecell{Compares\\verify\\policies} & \makecell{Faults\\matched\\across\\policies} & \makecell{Non-LLM\\oracle} & \makecell{Cost\\account-\\ing} \\
\midrule
ReliabilityBench \cite{gupta2026reliabilitybench} & \yes & \pt & \no & \no & \yes & \yes \\
ToolMaze \cite{zhu2026toolmaze} & \yes & \pt & \no & \no & \pt & \yes \\
Loud/Quiet \cite{kraishan2026loud} & \yes & \pt & \no & \no & \pt & \no \\
SAG \cite{park2026sag} & \no & \pt & \yes & \no & \yes & \yes \\
Sherlock \cite{ro2025sherlock} & \pt & \no & \yes & \no & \pt & \yes \\
Reconsider \cite{chen2026reconsider} & \no & \yes & \pt & \no & \yes & \pt \\
Verified calls \cite{mansoor2026verifiedtools} & \yes & \yes & \pt & \no & \pt & \no \\
\textbf{This work} & \yes & \yes & \yes & \yes & \yes & \yes \\
\bottomrule
\end{tabular}
\end{table}

% ----------------------------------------------------------------------------
\section{Problem Formulation}
\label{sec:problem}

An episode runs for at most $T$ steps. At step $t$ the agent, given goal $g$ and history
$h_t$, emits a tool call $a_t$. The environment produces a true result $r_t$ and
updates hidden state $s_{t+1}$; the agent instead observes $o_t = F_\sigma(a_t, r_t)$,
where $F_\sigma$ is a seeded fault channel that returns $r_t$ unchanged or a corrupted version
of it. A verification policy $\pi_v(a_t, o_t, h_t) \in \{0,1\}$ decides whether to invoke a
verifier, which may trigger recovery (re-calling the tool or re-planning). The episode
succeeds if the oracle $\mathcal{O}(s_T) = 1$.

Let $c_v$ be the cost of one verification (tokens, latency, calls). For a policy class
$\Pi$, the goal is
\begin{equation}
\max_{\pi_v \in \Pi} \; \Pr[\mathcal{O}(s_T)=1]
\quad \text{s.t.} \quad
\mathbb{E}\Big[\tfrac{1}{T}\textstyle\sum_t \pi_v(a_t,o_t,h_t)\Big] \le \rho ,
\label{eq:objective}
\end{equation}
i.e., the best reliability at a verification rate of at most $\rho$. Sweeping $\rho$ traces
a reliability--overhead frontier. A necessary condition for verification to help on a
faulted step is that $\pi_v$ fires on it; we call the fraction of faulted steps on which a
policy fires its \emph{trigger recall}. Trigger recall upper-bounds the benefit a perfect
verifier could deliver under that policy, and it can be measured without any LLM.

% ----------------------------------------------------------------------------
\section{The CostAware Testbed}
\label{sec:testbed}

\subsection{Procurement Simulator and Tools}
The simulator models nine entity types: vendors (approval, suspension), components,
supplier quotes (price, expiry, status), projects with procurement rules, budgets
(allocated, spent, reserved), purchase orders (POs), approvals, invoices, and shipments.
Agents act through 17 tools, each tagged with a risk tier in a single registry that also
supplies JSON argument schemas for tool calling (Table~\ref{tab:tools}).

\begin{table}[t]
\caption{Tool registry by risk tier.}
\label{tab:tools}
\centering
\footnotesize
\begin{tabular}{@{}p{0.17\columnwidth}p{0.75\columnwidth}@{}}
\toprule
Tier & Tools \\
\midrule
Low (read) & search\_vendors, get\_vendor\_details, check\_vendor\_status, get\_quotes,
check\_quote\_validity, check\_budget, check\_existing\_orders, compare\_quotes, check\_invoice \\
Medium & request\_discount, request\_approval, dispute\_invoice \\
High & create\_purchase\_order, modify\_order, approve\_order, release\_order, cancel\_order \\
\bottomrule
\end{tabular}
\end{table}

Two design choices make faults consequential. First, the simulator enforces \emph{workflow}
rules (unknown identifiers, illegal state transitions, releasing an order above the
approval threshold without a valid approval) but does not block \emph{policy} mistakes such
as ordering from a suspended vendor, on an expired quote, beyond budget, or twice. Those
actions succeed and are failed later by the oracle; if the simulator blocked them, a misled
agent could never fail and verification policies would be indistinguishable. Second, POs
are priced from quotes and identifiers are sequential, so reruns are bit-identical.
Modifying an order invalidates its approval, creating long-range dependencies.

\subsection{Tasks and Oracle}
The 20 tasks span 11 families (e.g., vendor status, quote conflict, budget constraint,
discount negotiation, approval workflow, PO modification and cancellation, invoice
mismatch, duplicate prevention): 7 easy, 8 medium, and 5 hard. Reference solutions range
from 2 to 14 tool calls (139 in total). Each task specifies a goal, an initial state, an
expected final state, a reference solution that must pass, and at least one negative
trajectory that must fail for the intended reason.

The oracle is pure Python and is tested to import no LLM library. It applies always-on
invariants (budgets within limits, every active order placed with a vendor in good standing,
quotes valid at order time, no duplicate active orders) and per-task checks on required
orders (vendor, quantity, status, price, valid approval), invoice states, and budget
reservations. Orders are located by project and component rather than by PO identifier.

\subsection{Fault Model and Matched Faults}
\label{sec:faults}
The injector applies one of eight fault types (Table~\ref{tab:faults}) to a tool's
observed result with probability equal to the fault rate. Faults reach nested fields (e.g.,
one quote inside a list) and a fault that would not change the result is skipped, so every
logged fault is a visible corruption. A timeout on a write tool still performs the write,
so a naive retry can create a duplicate order---the non-atomic failure studied in
\cite{mansoor2026verifiedtools}.

\begin{table}[t]
\caption{Fault types and their observable signature.}
\label{tab:faults}
\centering
\footnotesize
\begin{tabular}{@{}p{0.26\columnwidth}p{0.66\columnwidth}@{}}
\toprule
Fault & Effect on observed result \\
\midrule
wrong\_value & numeric field off by 10--30\% \\
stale\_info & older view of the same record \\
contradictory & boolean flag flipped (e.g., suspended) \\
malformed\_output & object replaced by a truncated raw string \\
missing\_field & a required key removed \\
timeout & no answer (write may still have happened) \\
incorrect\_status & status replaced by a plausible wrong one \\
corrupted\_obs & numeric field zeroed \\
\bottomrule
\end{tabular}
\end{table}

\emph{Matched faults.} A shared random stream would make faults depend on how many calls
preceded them, so a policy that verifies more (and therefore calls more tools) would face
different faults. Instead, each call's fault is derived from
$\mathrm{SHA256}(\sigma \,\|\, \text{tool} \,\|\, \text{canonical args} \,\|\, k)$, where $k$ is the
occurrence count of that exact call. The same agent action therefore receives the same fault
under every policy, while a recovery re-call is an independent draw. A replay module
re-executes a logged trajectory and reports any divergence from the recorded observations.

\subsection{Agent}
The agent is a LangGraph \cite{langgraph} state machine:
\textit{init} $\rightarrow$ \textit{agent} $\rightarrow$ \textit{tool\_exec} $\rightarrow$
\textit{fault\_injection} $\rightarrow$ \textit{policy\_router} $\rightarrow$ \textit{observe}
$\rightarrow$ (\textit{agent} $|$ \textit{finalize}). It issues one structured tool call per turn
in the ReAct style \cite{yao2023react}, ends via a \texttt{finish} tool, a step limit (30), or a
recursion backstop, and sees only observed results; true results reach only the logger and
oracle. Every step is logged as JSONL with true and observed results, the injected fault, and
the policy decision with its trigger reason.

\subsection{Verification Policies}
\begin{itemize}
\item \textbf{A (never)}: lower bound on cost and reliability.
\item \textbf{B (always)}: verify after every step.
\item \textbf{C (fixed schedule)}: verify every third step.
\item \textbf{D (adaptive)}: verify if the tool is high-risk (tool-based trigger), or the
observation fails a deterministic schema/sanity check---error or timeout, unparsable output,
missing budget fields, negative remaining budget, a vendor both approved and suspended, or a
PO creation without an identifier (content-based trigger). A third, confidence-based trigger
is specified but not yet wired to a model signal (Section~\ref{sec:limits}).
\item \textbf{E (budget-controlled)}: a continuous risk score compared to a threshold that an
EWMA controller adjusts toward a target verification rate $\rho$.
\end{itemize}

% ----------------------------------------------------------------------------
\section{Trigger-Coverage Analysis}
\label{sec:coverage}

\subsection{Setup}
To separate \emph{when} a policy verifies from \emph{how well} a verifier judges, we replay
each task's reference solution open-loop under fault injection and record, for every step,
whether each policy would trigger. Because faults are matched, all four policies see
identical faults. We use 50 seeds per task at fault rate 0.25, giving 6{,}950 tool calls of
which 1{,}764 (25.4\%) are faulted. Policy D's triggers are attributed to the tool-based or
content-based rule that fired first. No LLM is involved; the analysis is reproduced by
\texttt{scripts/trigger\_coverage.py}.

\subsection{Results}

\begin{table}[t]
\caption{Verification rate and trigger recall (rate 0.25, 6{,}950 steps).}
\label{tab:overall}
\centering
\footnotesize
\begin{tabular}{@{}lccc@{}}
\toprule
Policy & Verify rate & Recall (faulted) & Verify rate (clean) \\
\midrule
A (never) & 0.0\% & 0.0\% & 0.0\% \\
B (always) & 100.0\% & 100.0\% & 100.0\% \\
C (every 3rd) & 38.1\% & 39.8\% & 37.6\% \\
D (adaptive) & 43.7\% & 63.1\% & 37.1\% \\
\bottomrule
\end{tabular}
\end{table}

\begin{table}[t]
\caption{Trigger recall by fault type (rate 0.25). ``D content'' counts only D's
content-based trigger; the remainder of D's recall comes from the high-risk tool rule.}
\label{tab:perfault}
\centering
\footnotesize
\setlength{\tabcolsep}{4pt}
\begin{tabular}{@{}llrrrr@{}}
\toprule
Group & Fault & $n$ & C & D & D content \\
\midrule
Loud    & malformed\_output & 330 & 41.2 & 100.0 & 61.8 \\
        & timeout           & 312 & 37.5 & 100.0 & 75.0 \\
\midrule
Partial & missing\_field    & 376 & 33.2 & 42.3 & 8.0 \\
        & contradictory     & 135 & 39.3 & 40.7 & 13.3 \\
\midrule
Silent  & stale\_info       & 161 & 41.6 & 68.3 & 0.0 \\
        & incorrect\_status & 168 & 41.7 & 47.0 & 0.0 \\
        & wrong\_value      & 139 & 46.8 & 30.9 & 0.0 \\
        & corrupted\_obs    & 143 & 48.3 & 17.5 & 0.0 \\
\midrule
\multicolumn{2}{@{}l}{Silent (pooled)} & 611 & 44.4 & 42.1 & 0.0 \\
\multicolumn{2}{@{}l}{Loud (pooled)}   & 642 & 39.4 & 100.0 & 68.2 \\
\bottomrule
\end{tabular}
\end{table}

Table~\ref{tab:overall} shows the headline comparison. At a verification rate only 5.6 points
above the fixed schedule, Policy D reaches 63.1\% of faulted steps versus 39.8\%, while its
rate on clean steps (37.1\%) is no higher than C's. Taken alone, this would suggest D is
a strictly better allocation of verification budget.

Table~\ref{tab:perfault} shows where that advantage comes from. D catches every
\emph{loud} fault---timeouts and malformed output surface as errors or parse failures---and
these two types account for 642 of the 1{,}764 faulted steps. On the four \emph{silent} fault
types, which return well-formed results with wrong content, D's content trigger fires
\emph{zero} times. D still verifies some of these steps, but only when the faulted call
happens to be a high-risk tool; pooled over silent faults its recall (42.1\%) is no better
than the fixed schedule (44.4\%). On \texttt{wrong\_value} and \texttt{corrupted\_obs} it is
substantially worse (30.9\% vs.\ 46.8\%, 17.5\% vs.\ 48.3\%), because these faults mostly hit
low-risk read tools such as \texttt{check\_budget} and \texttt{get\_quotes}, which D never
verifies unless their output is malformed. Across all faults, only 27.6\% of faulted steps are
noticed by D's content trigger at all.

The pattern is stable across fault rates: at 0.10 and 0.40, D's recall is 61.6\% and 64.9\%,
its content trigger again fires on 0\% of silent faults, and the fixed schedule again exceeds
D on \texttt{wrong\_value} and \texttt{corrupted\_obs}.

\subsection{Interpretation}
The finding mirrors the agent-side result of Kraishan \cite{kraishan2026loud}---agents react to
the error channel, not to content---one level up: \emph{policies} built from risk tiers and
schema checks also react to the error channel. A silent fault on a read tool (a mis-reported
budget, a stale quote) is the input that leads to a wrong high-risk write two steps later,
and the risk-tier rule verifies the write but not the read that corrupted it. Whether
verifying the write is sufficient depends on the verifier re-querying the corrupted fact,
which is an empirical question for the end-to-end experiments. The analysis suggests two
design implications: (i) adaptive policies need a content-sensitive signal for in-format
faults, e.g., cross-observation consistency checks or model uncertainty; and (ii) risk should
be attributed to the \emph{reads that feed} a high-risk write, not only to the write.

% ----------------------------------------------------------------------------
\section{Evaluation Protocol}
\label{sec:protocol}

The following experiments are specified and supported by the harness but have not yet been
run; we state them so that hypotheses are fixed before results are seen.

\textbf{Design.} One agent model, policies A--E, 20 tasks, matched seeds, fault rate 0.25
with a sweep over $\{0.10, 0.25, 0.40\}$, and a clean (rate 0) control. The verifier and the
recovery protocol (re-call once, then re-plan with the flagged concern, then record an
unrecovered fault) are identical across policies; only the decision of \emph{when} to verify
changes.

\textbf{Metrics.} Task success (oracle pass rate); fault recovery rate (faulted episodes that
pass); error propagation (faulted episodes that fail); verification overhead (verifier calls per
action); tokens, dollar cost, and latency per task; and the reliability--overhead frontier of
Eq.~\eqref{eq:objective}. Verifier precision and recall are measured against the logged
ground-truth fault labels.

\textbf{Hypotheses.}
H1: Policy A's failure rate grows with horizon length and fault rate.
H2: Policy B improves reliability over A at the highest overhead.
H3: Policy D approaches B's reliability at lower overhead.
H4: The high-risk and schema triggers each contribute to D's reliability (signal ablation).
H5 (from Section~\ref{sec:coverage}): D's reliability advantage over C is concentrated on loud
faults and vanishes or reverses on silent faults; adding a content-sensitive trigger closes
this gap.

\textbf{Controls.} A random-verification policy matched to D's per-task verification rate
separates the value of D's \emph{signals} from the value of its \emph{budget}. Results will be
reported with 95\% bootstrap confidence intervals over seeds and paired tests across policies,
which matched faults make valid.

% ----------------------------------------------------------------------------
\section{Limitations and Threats to Validity}
\label{sec:limits}

\textbf{Open-loop replay.} Section~\ref{sec:coverage} replays fixed reference solutions; an LLM
agent will take different paths, and its own reactions to faults (e.g., re-querying) are not
captured. Trigger recall bounds, but does not measure, the reliability benefit of
verification.

\textbf{Triggering is not detection.} A triggered verifier must still recognize the fault.
Self-verification without external evidence is unreliable \cite{huang2024cannot}, so end-to-end
benefit may be well below trigger recall.

\textbf{Confidence trigger not implemented.} D's specified low-confidence trigger reads a
confidence field that no tool or model currently supplies, so the D evaluated here is the
two-signal (risk + schema) variant. A logprob- or entropy-based signal \cite{park2026sag} may
recover some silent-fault coverage; this is part of H5.

\textbf{Synthetic environment and fault mix.} The world is synthetic and small (20 tasks). The
effective mix over-represents faults that always apply (malformed, missing field, timeout)
relative to those requiring a particular field. Per-type results are reported to make this
visible; the pooled numbers should be read with it in mind.

\textbf{Relation to prior work.} The literature on this topic is moving quickly (several of the
works cited appeared in the last three months); positioning in Table~\ref{tab:positioning} is
based on published abstracts and should be re-checked against full texts.

% ----------------------------------------------------------------------------
\section{Conclusion}
CostAware isolates the question of \emph{when} a tool-using agent should verify, by holding
faults fixed across verification policies in a stateful environment judged by a deterministic
oracle. An LLM-free analysis already yields a concrete result: a risk-tier and schema-based
adaptive policy looks markedly more efficient than a fixed schedule, but its advantage comes
entirely from faults that announce themselves; on silent, in-format faults it is no better
than verifying every third step. The planned end-to-end experiments test whether this gap
persists once an LLM verifier and recovery are in the loop, and whether content-sensitive
signals close it.

% ----------------------------------------------------------------------------
\section*{Code and Data Availability}
Simulator, tasks, injector, oracle, agent, and analysis scripts:
\url{https://github.com/SP-00-Origin/costAware}.

\section*{Acknowledgment}
% IEEE requires disclosing AI-generated content (text, figures, code) in the acknowledgments,
% naming the system and the sections it was used for. Edit to reflect actual use.
The authors used an AI assistant (Claude, Anthropic) to help survey related work and draft
portions of this manuscript; all content was reviewed and verified by the authors, who take
full responsibility for it.

\bibliographystyle{IEEEtran}
\bibliography{references}

\end{document}
